\documentclass{article}
\usepackage{amsmath}
\usepackage{siunitx}
\usepackage{bm} % For bold math symbols


\usepackage{geometry}% 1-inch margins
\geometry{letterpaper, margin=1in}% 1-inch margins


\title{EMCH 290 HW\#2 Solutions}
\date{Fall 2023}
\author{Jacob C. Vaught}
\begin{document}

\maketitle

\newpage
\section*{Question 1 (7 Points)}
A flywheel whose moment of inertia \( I = \SI{200}{kg \cdot m^2} \) is spinning at the speed shown in the figure. The flywheel is connected to a mass via a pulley. For a flywheel, 
\[
KE = \frac{I \omega^2}{2},
\]
where \( \omega \) is the angular velocity in \si{rad/s}. Let \( g = \SI{9.81}{m/s^2} \). Assume air resistance and friction are negligible.

\subsection*{Part 1: \( m = \SI{50}{kg} \)}
To find the change in height of the mass in meters, equate kinetic energy (KE) to potential energy (PE):
\[
KE = \frac{I \omega^2}{2} = \frac{200 \, \text{kg} \cdot \text{m}^2 \times \left( \frac{200 \, \text{RPM} \times 2\pi \, \text{rad}}{60 \, \text{s}} \right)^2}{2} = \SI{43864}{J} = \SI{43.864}{kJ}
\]
\[
PE = mgh = KE \implies mgh = \SI{43864}{J} \implies h = \frac{43864}{50 \times 9.81} = \SI{89.42}{m}
\]

\subsection*{Part 2: \( m = \SI{140}{kg} \)}
To find the change in height when the mass is \SI{140}{kg}, again equate KE to PE:
\[
PE = mgh = KE \implies mgh = \SI{43864}{J} \implies h = \frac{43864}{140 \times 9.81} = \SI{31.94}{m}
\]



\newpage
\section*{Question 2 (8 Points)}

A foreign country is aiming a missile at a target destination. However, the missile is accidentally fired straight upward and reaches a height of \SI{900}{m} when its velocity reaches zero (apogee). The only force acting on the object is the force of gravity (assume air resistance is negligible). \( g = \SI{9.81}{m/s^2} \).

\subsection*{Part 1: Kinetic Energy Needed to Fire the Missile}
The mass \( m \) of the missile is \( \SI{90}{kg} \). To find the kinetic energy needed to fire the missile, we equate potential energy (PE) to kinetic energy (KE) at the apogee:
\[
KE = \frac{1}{2} m v^2
\]
\[
PE = mgh
\]
Therefore, the kinetic energy required to reach this height is:
\[
KE = mgh_f = 90 \times 900 \times 9.81 \Rightarrow \mathbf{\SI{794610}{J} = \SI{794.610}{kJ}}
\]

\subsection*{Part 2: Firing Velocity of the Missile}
To find the initial velocity \( v_i \) of the missile, we use the equation:
\[
\frac{1}{2} m v_i^2 = mgh_f \Rightarrow v_i = \sqrt{2gh} = \sqrt{2 \times 9.81 \times 900} \Rightarrow \mathbf{\SI{133.365}{m/s} = \SI{0.133365}{km/s}}
\]

\newpage
\section*{Question 3 (9 Points)}

A vehicle is stopped at the red light on Main St. and Blossom St. While waiting, the vehicle's brakes fail, causing the vehicle to roll down the hill towards Swearingen. Assuming the street is inclined at \(60^\circ\) from the horizontal, air resistance and friction between the vehicle and the street are negligible. Let \( g = \SI{9.81}{m/s^2} \).

\subsection*{Part 1: Kinetic Energy at the Bottom of the Street}
To find the kinetic energy of the vehicle at the bottom of the hill, we can equate the kinetic energy (KE) to the gravitational potential energy (PE):
\[
KE = \frac{1}{2} mv^2
\]
\[
PE = mgh
\]
The energy before rolling down is converted entirely into kinetic energy:
\[
mgh_f = \frac{1}{2} mv_f^2 \Rightarrow KE = mgh_f = 2000 \times 9.81 \times 100 \times \sin(60^\circ) \Rightarrow \mathbf{\SI{1699100.4}{J} = \SI{1699.1}{kJ}}
\]

\subsection*{Part 2: Velocity at the Bottom of the Street in m/s}
To find the velocity \( v_f \) of the vehicle at the bottom, we can rearrange the equation:
\[
v_f = \sqrt{2 \times gh_f} = \sqrt{2 \times 9.81 \times 100 \times \sin(60^\circ)} \Rightarrow \mathbf{\SI{41.22}{m/s}}
\]

\subsection*{Part 3: Velocity at the Bottom of the Street in mph}
To convert the velocity from \( m/s \) to \( \text{mph} \), we use the conversion factor \( 1 \, \text{m/s} = 2.23694 \, \text{mph} \):
\[
v_f = 41.22 \times 2.23694 \Rightarrow \mathbf{\SI{92.14}{mph}}
\]

\newpage
\section*{Question 4 (7 Points)}

You're enjoying a day at the beach when you spot a sailboat, weighing \SI{500}{kg}, with a \SI{100}{kg} sailor aboard. The boat is making good speed when, out of nowhere, a storm brews. The storm's force, a formidable \SI{1000}{N}, causes the sailboat to come to a rest, before it capsizes the sailboat. Assume the force the wind applies is parallel to the direction of the boat initially.

\subsection*{Part 1: Energy Transfer by Work}

To find the work done by the force of the storm, we use the work-energy principle:
\[
W = \frac{1}{2}m(v_f^2 - v_i^2)
\]
Here, \( m = \SI{500}{kg} + \SI{100}{kg} = \SI{600}{kg} \) (mass of boat and sailor combined), \( v_i = \SI{40}{m/s} \), and \( v_f = \SI{0}{m/s} \).
\[
W = \frac{1}{2} \times 600 \times (0 - 40^2) = -\SI{480000}{J} \Rightarrow \mathbf{-\SI{480}{kJ}}
\]

\subsection*{Part 2: Distance Travelled}

To find the distance \( \Delta x \) the sailboat managed to travel before the storm overpowered it, we use the equation \( W = F \Delta x \):
\[
\Delta x = \frac{W}{F} = \frac{-\SI{480000}{J}}{-\SI{1000}{N}} \Rightarrow \mathbf{\SI{480}{m}}
\]

\newpage
\section*{Question 5 (10 Points)}

A Formula One car tire after a race has \SI{10}{kg} of air in it with a pressure of \( p_1 = \SI{450}{kPa} \), and a specific volume of \( v_1 = \SI{0.25}{m^3/kg} \). After the tire sits for a while, the tire's pressure and specific volume change to \( p_2 = \SI{140}{kPa} \), \( v_2 = \SI{0.7}{m^3/kg} \).

\subsection*{Part 1: Polytropic Exponent \( n \) for the Process}

We start with the equation for a polytropic process, defined as \( pV^n = \text{constant} \) or \( pv^n = \text{constant} \), and solve for \( n \):
\[
n = \frac{\ln\left(\frac{p_1}{p_2}\right)}{\ln\left(\frac{v_2}{v_1}\right)} = \frac{\ln\left(\frac{450}{140}\right)}{\ln\left(\frac{0.7}{0.25}\right)} \Rightarrow \mathbf{1.134}
\]

\subsection*{Part 2: Energy Transfer by Work}

The work done \( W \) during a polytropic process is given by the integral \(\int PdV\), which can be expressed as:
\[
W = m \left( \frac{p_2v_2 - p_1v_1}{1-n} \right)
\]
Here, \( m = \SI{10}{kg} \), \( p_1 = \SI{450000}{Pa} \), \( p_2 = \SI{140000}{Pa} \), \( v_1 = \SI{0.25}{m^3/kg} \), \( v_2 = \SI{0.7}{m^3/kg} \), and \( n = 1.134 \).
\[
W = 10 \left( \frac{140000 \times 0.7 - 450000 \times 0.25}{1-1.134} \right) \Rightarrow \mathbf{\SI{1082089}{J} = \SI{1082}{kJ} = \SI{1.08}{MJ}}
\]

\newpage
\section*{Question 6 (7 Points)}

A heat pump cycle operating at steady state receives energy by heat transfer from well water at \SI{10}{\degreeCelsius} and discharges energy by heat transfer to a building at the rate of \(\SI{1.2e5}{kJ/h}\). Over a period of 14 days, an electric meter records that \SI{1620}{kW \cdot h} of electricity is provided to the heat pump. These are the only energy transfers involved.

\subsection*{Part 1: Total Energy Transfer by Work}

The total amount of energy transfer by work into the heat pump over the 14-day period can be calculated as follows:
\[
1620\left[\text{kW}\cdot \text{h}\right]\times\left[\frac{3600\text{s}}{\text{h}}\right]\times\left[\frac{\text{kJ}}{\text{kW}\cdot \text{s}}\right]\times\left[\frac{{10}^{-6}\text{GJ}}{\text{kJ}}\right]=\mathbf{5.832}\left[\text{GJ}\right]
\]

\subsection*{Part 2: Energy from Well Water}

The amount of energy that the heat pump receives by heat transfer from the well water over the 14-day period is:
\[
Q_{\text{out}} = 1.2\times{10}^5\left[\frac{\text{kJ}}{\text{h}}\right]\times 14\left[\text{days}\right]\times\left[\frac{24\text{h}}{\text{day}}\right]\times\left[\frac{{10}^{-6}\text{GJ}}{\text{kJ}}\right] = \mathbf{40.32}\left[\text{GJ}\right]
\]

\subsection*{Part 3: Coefficient of Performance}

The coefficient of performance \( \gamma \) can be calculated as:
\[
\gamma = \frac{Q_{\text{out}}}{W_{\text{cycle}}} = \frac{40.32}{5.832} = \mathbf{6.9135}
\]

\newpage
\section*{Question 7 (7 Points)}

Imagine you're working as an engineer in a high-tech aerospace company, and you're currently testing a new model of an electric motor for a future drone design. The electric motor, while operating under steady-state conditions, draws a current of \(11 \, \text{A}\) at a voltage of \(120 \, \text{V}\). During operation, the output shaft of the motor develops a torque of \(8.4 \, \text{N} \cdot \text{m}\) and rotates at a speed of \(1200 \, \text{RPM}\).

\[
\dot{W} = P = V \times I
\]

\subsection*{Part 1: Electric Power}

The electric power drawn by the motor can be calculated as follows:

\[
\dot{W} = -E \times I = -120 \, \text{V} \times 11 \, \text{A} = -\mathbf{1320} \, \text{W} = -\mathbf{1.32} \, \text{kW}
\]

\subsection*{Part 2: Power Developed by the Output Shaft}

The power developed by the output shaft is given by:

\[
\dot{W}_{\text{shaft}} = 8.4 \, \text{N} \cdot \text{m} \times 1200 \, \text{RPM} \times \left[\frac{2\pi}{\text{rad}}\right] \times \left[\frac{1 \, \text{min}}{60 \, \text{s}}\right] = \mathbf{1055.6} \, \text{W} = \mathbf{1.055} \, \text{kW}
\]

\subsection*{Part 3: Rate of Heat Transfer}

The rate of heat transfer in the system can be calculated as:

\[
\dot{Q}_{\text{system}} = 1.055 - 1.32 = -\mathbf{0.265} \, \text{kW}
\]

\newpage
\section*{Question 8 (7 Points)}

An object whose mass is \(100 \, \text{kg}\) experiences a decrease in kinetic energy of \(500 \, \text{N} \cdot \text{m}\) and an increase in potential energy of \(1500 \, \text{N} \cdot \text{m}\). The initial velocity and elevation of the object, each relative to the surface of the Earth, are \(40 \, \text{m/s}\) and \(30 \, \text{m}\), respectively. Given that \(g = 9.81 \, \text{m/s}^2\), determine:

\subsection*{Part 1: Final Velocity}

The initial kinetic energy \(KE_{\text{initial}}\) is:

\[
KE_{\text{initial}} = \frac{1}{2} mv^2 = \frac{1}{2} \times 100 \times 40^2 = 80000 \, \text{N} \cdot \text{m}
\]
The final kinetic energy \(KE_{\text{final}}\) becomes:

\[
KE_{\text{final}} = 80000 - 500 = 79500 \, \text{N} \cdot \text{m}
\]
The final velocity \(V_{\text{final}}\) is:

\[
V_{\text{final}} = \sqrt{\frac{KE_{\text{final}}}{0.5m}} = \sqrt{\frac{79500}{0.5 \times 100}} = \mathbf{39.874} \, \text{m/s}
\]

\subsection*{Part 2: Final Elevation}

The initial potential energy \(PE_{\text{initial}}\) is:

\[
PE_{\text{initial}} = mgh = 100 \times 9.81 \times 30 = 29430 \, \text{N} \cdot \text{m}
\]
The final potential energy \(PE_{\text{final}}\) is:

\[
PE_{\text{final}} = 29430 + 1500 = 30930 \, \text{N} \cdot \text{m}
\]
The final elevation \(h_{\text{final}}\) is:

\[
h_{\text{final}} = \frac{PE_{\text{final}}}{mg} = \frac{30930}{100 \times 9.81} = \mathbf{31.53} \, \text{m}
\]

\newpage
\section*{Question 9 (7 Points)}

Find the mass flow rate (in grams per second) of water for a control volume system, if the volume flow rate is \(2 \, \text{liter per second}\). \(\rho_{\text{water}}=1 \, \left[\frac{g}{\text{cm}^3}\right]\).
\newline
First, convert the volume flow rate from liters to cubic centimeters:
\[
2 \, \left[\text{Liter}\right] = 2000 \, \left[\text{cm}^3\right]
\]
\newline
The mass flow rate is then calculated as:
\[
\text{Mass flow rate} = 2000 \, \left[\frac{\text{cm}^3}{\text{s}}\right] \times 1 \, \left[\frac{g}{\text{cm}^3}\right] = \mathbf{2000} \, \left[\frac{\text{g}}{\text{s}}\right]
\]
\newpage
\section*{Question 10 (7 Points)}

What is the change in the internal energy of the system if the surrounding releases \(300 \, \text{J}\) of heat energy and if the system does \(550 \, \text{J}\) of work on the surroundings?
\newline
The change in internal energy is given by:
\[
\text{Change in internal energy} = \left(+300\right) + \left(-550\right) = -\mathbf{250} \, \left[\text{J}\right]
\]

\newpage
\section*{Question 11 (7 Points)}

The refrigerator shown in the figure steadily receives a power input of \(0.15 \, \text{[kW]}\) while rejecting energy by heat transfer to the surroundings at a rate of \(0.6 \, \text{[kW]}\). Determine the rate at which energy is removed by heat transfer from the refrigerated space, in \([kW]\), and the refrigerator's coefficient of performance.
\newline \newline
To find the rate of energy removal from the refrigerated space and the coefficient of performance, we use the following equations:
 \newline
\[
W_{\text{cycle}} = Q_{\text{out}} - Q_{\text{in}}; \quad \beta = \frac{Q_{\text{in}}}{W_{\text{cycle}}}
\]
 \newline
Substituting the given values:
 \newline
\[
Q_{\text{in}} = Q_{\text{out}} - W_{\text{cycle}} = 0.6 - 0.15 = \mathbf{0.45} \, \text{[kW]}
\]
\[
\beta = \frac{0.45}{0.15} = \mathbf{3}
\]
 \newline
Thus, the rate at which energy is removed from the refrigerated space is \(0.45 \, \text{[kW]}\) and the coefficient of performance is \(3\).

\newpage
\section*{Question 12 (7 Points)}

For a power cycle operating as shown in the figure, \( W_{\text{cycle}} = 600 \, \text{kJ} \) and \( Q_{\text{out}} = 1800 \, \text{kJ} \).

\subsection*{Part 1: Heat Input}

The heat input \( Q_{\text{in}} \) can be calculated using the equation:

\[
Q_{\text{in}} = Q_{\text{out}} + W_{\text{cycle}} = 1800 + 600 = \mathbf{2400} \, \text{kJ}
\]

\subsection*{Part 2: Thermal Efficiency}

The thermal efficiency \( \eta \) of the cycle can be calculated as:

\[
\eta = \frac{W_{\text{cycle}}}{Q_{\text{in}}} = \frac{600}{2400} = \mathbf{0.25}
\]

Thus, the heat input \( Q_{\text{in}} \) is \( 2400 \, \text{kJ} \) and the thermal efficiency of the cycle is \( 0.25 \).

\newpage
\section*{Question 13 (10 Points)}

A vertical piston–cylinder assembly with a piston of mass 35 kg and having a face area of \(0.006 \, \text{m}^2\) contains air. The mass of air is \(0.0025 \, \text{kg}\), and initially the air occupies a volume of \(0.0025 \, \text{m}^3\). The atmosphere exerts a pressure of 100 kPa on the top of the piston. The volume of the air slowly decreases to \(0.001 \, \text{m}^3\) as energy with a magnitude of 1 kJ is slowly removed by heat transfer. Friction between the piston and the cylinder wall can be neglected and \(g = 9.8 \, \text{m/s}^2\).

\subsection*{Part 1: Work Done}

The sum of forces in the vertical direction is given by:
\[
\sum F_z = p_{\text{air}} \times A_{\text{piston}} - p_{\text{atm}} \times A_{\text{piston}} - m_{\text{piston}} \times g = 0
\]
Solving for \(p_{\text{air}}\):
\[
p_{\text{air}} = p_{\text{atm}} + \frac{m_{\text{piston}} \times g}{A_{\text{piston}}} = 100000 + \frac{35 \times 9.81}{0.006} = 157225 \, \text{Pa} = 157.22 \, \text{kPa}
\]
The work \(W\) done by the air can be calculated as:
\[
W = \int_{V_1}^{V_2}{p \, dV} = p \times (V_2 - V_1) = 157225 \times (0.001 - 0.0025) = -235.8375 \, \text{J} = -\mathbf{0.2358} \, \text{kJ}
\]

\subsection*{Part 2: Change in Specific Internal Energy}

The change in internal energy \( \Delta U \) is given by:
\[
\Delta U = Q - W = -1 - (-0.2358) = -0.7642 \, \text{kJ}
\]
The change in specific internal energy \( \Delta u \) can be calculated as:
\[
\Delta u = \frac{\Delta U}{m} = -\frac{0.7642}{0.0025} = -\mathbf{305.68} \, \frac{\text{kJ}}{\text{kg}}
\]

\end{document}
